\section{システムの動作}
\subsection{動作確認の手順}
3部門を結合して動作確認を行った．
具体的には，最初にシリアルモニタに"chiba"と入力し，光通信→音通信→振動通信の順で通信を行っていき，各部門にかかる時間と受け取った4進数または2進数の数列，最終的にデコードされた文字を記録した．
動作確認時の条件は以下の通りである．
\begin{itemize}
    \item 各Arduinoに繋いだ電源はすべてMac bookのUSB type-Cからのみの給電
    \item スピーカとマイクの距離は2 cm
    \item 振動モータと加速度センサは，4つのゴムで脚をつくった薄いアクリル板に貼り付け，その距離は2 cm
    \item 時間の計測はiphone 14の「時計」アプリのラップタイムを利用
\end{itemize}

\subsection{目標の設定}
本システムが有効であるかの指標として，各部門に表\ref{table:sihyo}のような目標を設定する．
これらの目標に到達することで，有効なシステムであると判断する．
bps（bit par second）は，1秒あたりに送信されたビット数を表す．
フレーム成功率（FSR）は，送信されたフレームの総数に対する誤りがなく通信できたフレームの割合で，総フレーム数を$N_a$，誤り無く送信されたフレーム数を$N_c$として式\ref{eq:fsr}で表される．
ビット誤り率（BER）は，送信されたビット数に対する誤って受信したビット数の割合で，総ビット数を$N_t$，誤りビット数を$N_e$として，式\ref{eq:ber}で表される．
\begin{gather}
    FSR = \frac{N_c}{N_a} \label{eq:fsr} \\
    BER = \frac{N_t}{N_e} \label{eq:ber}
\end{gather}

光通信において，FSR $\geq$ 0.90，100 bps以上であるとシステムが有効であるといえる\cite{ref:fsr}\cite{ref:lifi}．
また，データを音のオンオフで符号化した音響通信の例では，雑音が通信音量以下だとBERは約2 \%，通信音量以上だとBERが10 \%を超えたという報告がある\cite{ref:ber}．
この研究\cite{ref:ber}における通信速度は 10bpsであり，この条件でBERが2$\simeq$10 \%であるため，10 bpsで適切な通信が可能だと考える．
これらを踏まえ，$FSR \geq 0.90$，$BER < 10{-1}$，100 bps以上（光通信），10 bps以上（音，振動通信）を目標として設定する．


\begin{table}[ht]
    \centering
    \caption{各部門における到達目標}
    \label{table:sihyo}
    \begin{tabular}{|c|c|c|} \hline
        部門 &通信精度  &通信速度  \\ \hline
        光通信 &$FSR \geq 0.90$ &100 bps以上 \\ \hline
        音通信 &$BER \geq 10^{-1}$ &10 bps以上 \\ \hline
        振動通信 &$BER \geq 10^{-1}$ &10 bps以上 \\ \hline
    \end{tabular}
\end{table}


\subsection{動作結果}
計12回の試行を行った結果を表\ref{table:result}に，各部門で変換した4進数または2進数を表\ref{table:result2}，表\ref{table:result3}に示す．
光，音，振動通信と全体の経過時間の平均はそれぞれ1.07 s，5.23 s，14.20 s，20.51 sだった．
光通信では一度も間違って通信することはなかったが，音通信では12回中2回，振動通信では12回中11回誤った情報が伝送された．
デコード結果は，完全な"chiba"が1回，1文字誤りが2回，2文字誤りが1回だった．


\begin{table}[htbp]
    \centering
    \caption{本システムの動作結果}
    \label{table:result}
    \begin{tabular}{|c|c|c|c|c|c|} \hline
        試行回数 &光通信の時間[s]  &音通信の時間[s] &振動通信の時間[s] &合計時間[s] &デコード結果  \\ \hline
        1 &1.06   &5.28  &13.73  &20.07 &gl\{bi  \\ \hline
        2 &1.09   &5.21  &12.92  &19.22 &chibq  \\ \hline
        3 &1.09   &5.19  &13.08  &19.36 &s|<s0  \\ \hline
        4 &1.08   &5.21  &20.92  &27.21 &chija8 \\ \hline
        5 &1.09   &5.04  &12.26  &18.39 &gH     \\ \hline
        6 &1.01   &5.25  &13.08  &19.34 &cBisA  \\ \hline
        7 &1.05   &5.13  &15.14  &21.32 &chkba  \\ \hline
        8 &1.56   &4.75  &16.50  &22.81 &oPIRx  \\ \hline
        9 &1.10   &6.08  &13.11  &20.29 &s|tPp  \\ \hline
        10 &0.98  &5.31  &12.98  &19.27 &chiba  \\ \hline
        11 &0.98  &5.21  &13.23  &19.42 &s|\{`q  \\ \hline
        12 &0.91  &5.06  &13.49  &19.46 &cv6\&  \\ \hline
    \end{tabular}
\end{table}

\begin{table}[htbp]
    \centering
    \caption{光，音通信の変換した4進数または2進数}
    \label{table:result2}
    \begin{tabular}{|c|c|c|c|} \hline
        試行回数  &デコード結果 &光受信（4進数） &音受信（2進数） \\ \hline
        1   &gl\{bi &12031220122112021201  &11000111101000110100111000101100001  \\ \hline
        2   &chibq  &12031220122112021201  &11000111101000110100111000101100001  \\ \hline
        3   &s|<s0  &12031220122112021201  &11000111101000110100111000101100001  \\ \hline
        4   &chija8 &12031220122112021201  &11000111101000110100111000101100001  \\ \hline
        5   &gH     &12031220122112021201  &1100011110100000101100100110         \\ \hline
        6   &cBisA  &12031220122112021201  &11000111101000110100111000101100001  \\ \hline
        7   &chkba  &12031220122112021201  &11000111101000110100111000101100001  \\ \hline
        8   &oPIRx  &12031220122112021201  &11000111101000110100111000101100001  \\ \hline
        9   &s|tPp  &12031220122112021201  &11000111101000110100111000101100001  \\ \hline
        10  &chiba  &12031220122112021201  &11000111101000110100111000101100001  \\ \hline
        11  &s|\{`q &12031220122112021201  &11000111101000110100111000101100001  \\ \hline
        12  &cv6\&  &12031220122112021201  &1100011110011000101100100110         \\ \hline
    \end{tabular}
\end{table}

\begin{table}[htbp]
    \centering
    \caption{振動通信の変換した4進数または2進数}
    \label{table:result3}
    \begin{tabular}{|c|c|p{10cm}|} \hline
        試行回数  &デコード結果 &振動受信（2進数） \\ \hline
        1   &gl\{bi &1100111111110011110111100010110100100000000000  \\ \hline
        2   &chibq  &記録ミス \\ \hline
        3   &s|<s0  &記録ミス  \\ \hline
        4   &chija8 &11000111101000110100111010101100001000000000010000111000 \newline
                     00000001000100000000000 \\ \hline
        5   &gH     &11001111001000001011001011100000000000     \\ \hline
        6   &cBisA  &1100011100001011010011110011100000100000000000  \\ \hline
        7   &chkba  &1100011110100011010111100010110000100100001100000000000  \\ \hline
        8   &oPIRx  &11011111010000100100110100101111000000001000100000000000  \\ \hline
        9   &s|tPp  &11100111111100111010010100001110000100000000000  \\ \hline
        10  &chiba  &1100011110100011010011100010110000100000000000  \\ \hline
        11  &s|\{`q &11100111111100111101111000001110001100000000000  \\ \hline
        12  &cv6\&  &11000111110110011011001001100000000100000000000  \\ \hline
    \end{tabular}
\end{table}

\clearpage